%%%%%%%%%%%%%%%%%%%%%%%%%%%%%%%%%%%%%%%%%%%%%%%%%%%%%%%%%%%%%%%%%%%%%%%%%%
%   FEHIME CEREN AY - CV
%   Tilpasset: Senior analytiker, brukerinnsikt og måling
%   Nav, Arbeids- og velferdsdirektoratet (FINN-kode 470275492)
%%%%%%%%%%%%%%%%%%%%%%%%%%%%%%%%%%%%%%%%%%%%%%%%%%%%%%%%%%%%%%%%%%%%%%%%%%

\documentclass[10pt,a4paper,roman]{moderncv}

\moderncvstyle{casual}
\moderncvcolor{blue}

\usepackage[utf8]{inputenc}
\usepackage[norsk]{babel}
\usepackage[scale=0.78,top=1.6cm,bottom=1.4cm]{geometry}
\setlength{\hintscolumnwidth}{3.4cm}
\usepackage{enumitem}
\setlist[itemize]{leftmargin=*,topsep=2pt,itemsep=1pt,parsep=0pt}

\newcommand{\cvbullets}[1]{\begin{itemize}#1\end{itemize}}

%----------------------------------------------------------------------------------------
%   NAVN OG KONTAKT
%----------------------------------------------------------------------------------------

\firstname{Fehime Ceren}
\familyname{Ay}

\title{Senioranalytiker \textbar{} brukerinnsikt, atferd og effektmåling}
\address{Bestumveien 37C, 0281}{Oslo}
\phone{+47 92 25 66 45}
\email{fcerenay@gmail.com}
\extrainfo{\href{https://cerenay.github.io}{\color{blue}cerenay.github.io}}

\begin{document}
\makecvtitle

%----------------------------------------------------------------------------------------
%   PROFIL
%----------------------------------------------------------------------------------------

\section{Profil}
\cvitem{}{
Analytiker med PhD i økonomi og fem års erfaring med å bruke data til å forstå
hvordan tjenester faktisk brukes og oppleves. I Skatteetaten jobber jeg med
brukeropplevelse i et tverrfaglig miljø med utviklere og designere, i
skjæringspunktet mellom atferd, analyse og effektmåling: jeg designer
spørreundersøkelser og eksperimenter, kombinerer registerdata med atferdsdata fra
Matomo og Skyra, og omsetter funnene til beslutningsgrunnlag som ledelse og
fagmiljøer kan handle på. Fagbakgrunnen min er atferdsøkonomi, og jeg er opptatt av
menneskene bak tallene: hvorfor folk gjør som de gjør, hva som hindrer dem, og
hva vi kan endre. Jeg er samtidig tydelig på hva data faktisk viser og hvor de
kommer til kort. Til daglig jobber jeg i SQL, Python og R på Databricks, og med
Matomo, Skyra og Splunk.
}

%----------------------------------------------------------------------------------------
%   NØKKELKOMPETANSE
%----------------------------------------------------------------------------------------

\section{Nøkkelkompetanse}

\cvitem{Brukerinnsikt \& måling}{Design av spørreundersøkelser og måleplaner,
toppoppgavemetodikk, måling av brukeratferd og brukerreiser på digitale flater.
Verktøy: Matomo, Skyra, Splunk}
\cvitem{Metode}{Effektmåling og kausal inferens, regresjonsanalyse,
eksperimentdesign og A/B-testing, segmentering, hypotesedrevet arbeid,
kvasi-eksperimentelle metoder}
\cvitem{Programmering}{SQL (uttrekk, modellering, optimalisering), Python, R}
\cvitem{Formidling}{Power BI, R Shiny og Quarto til dashbord, målekort og
selvbetjente analyseverktøy for team, ledelse og styringsmiljøer}
\cvitem{Plattform \& data}{Databricks, skybaserte og lokale datamiljøer,
Git/versjonskontroll, arbeid med store register-, survey- og hendelsesdatasett}
\cvitem{Datakvalitet}{Datavalidering og avviksanalyse, dokumentasjon av
definisjoner og måltall, etablering av felles datagrunnlag på tvers av kilder}
\cvitem{Språk}{Norsk (flytende), Engelsk (flytende), Tyrkisk (morsmål)}
\cvitem{Øvrig}{Norsk og tyrkisk statsborgerskap. Referanser oppgis på forespørsel.}

%----------------------------------------------------------------------------------------
%   ARBEIDSERFARING
%----------------------------------------------------------------------------------------

\section{Arbeidserfaring}

\cvitem{09/2023 -- nå}{\textbf{Seniorrådgiver, brukeropplevelse, atferd og effekt} \newline
\textit{Skatteetaten, Skattedirektoratet -- Data, Analyse og Kunnskap}}
\cvitem{}{\vspace{-4pt}
\cvbullets{
\item Arbeider med brukeropplevelse i tverrfaglige team sammen med utviklere og
      designere, og sørger for at endringer i tjenestene måles og følges opp.
\item Måler brukeropplevelse og bruksmønstre på digitale flater med Matomo, Skyra
      og Splunk, og kombinerer dette med register- og surveydata i Databricks.
\item Designer og analyserer spørreundersøkelser, eksperimenter og A/B-tester for
      å forstå brukeratferd og måle effekten av tiltak i etatens tjenester.
\item Kartlegger hva som faktisk hindrer folk i å bruke tjenestene riktig, blant
      annet digital sårbarhet, misforståelser av regelverk og betalingsproblemer,
      og bruker segmentering til å foreslå ulike tiltak til ulike grupper.
\item Leverer styringsinformasjon og beslutningsgrunnlag til ledelsen: målekort,
      nøkkeltall og rapporter som brukes i prioritering og oppfølging.
\item Bygger selvbetjente dashbord og analyseverktøy slik at fagmiljøene kan følge
      egne tall uten å være avhengige av én person.
\item Arbeider med datakvalitet og begrepsforvaltning, og samarbeider tett med IT-,
      analyse- og AI-miljøer. Deltar i internasjonale fagmiljøer (OECD, IBPPA).
}}

\cvitem{01/2021 -- 09/2023}{\textbf{Research Scientist / Dataanalytiker} \newline
\textit{Telenor ASA -- Research \& Innovation, Advanced Analytics \& AI}}
\cvitem{}{\vspace{-4pt}
\cvbullets{
\item Utviklet måling av brukertilfredshet for digitale kundeservicekanaler ved å
      analysere samtalelogger fra chatbot med prediktive modeller og tekstanalyse.
      Arbeidet er publisert i \textit{Quality and User Experience}.
\item Analyserte bruks- og aktivitetsdata fra digitale tjenester på tvers av
      forretningsenheter i Norden og Asia, og koblet interne data mot eksterne
      kilder for å svare på konkrete forretningsspørsmål.
\item Utformet spørreundersøkelser og eksperimenter, og strukturerte og
      kvalitetssikret rådata til analyseklare datasett som andre team kunne gjenbruke.
\item Formidlet funn og anbefalinger til produktledelse, teknologimiljøer og
      konsernledelse, tilpasset hver målgruppe.
\item Tilknyttet forsker ved NHH, FAIR Research Centre.
}}

\cvitem{08/2016 -- 12/2020}{\textbf{PhD-stipendiat} \newline
\textit{Norges Handelshøyskole (NHH) -- Institutt for økonomi, FAIR / The Choice Lab}}
\cvitem{}{\vspace{-4pt}
Eksperimentell og kvantitativ forskning på hvordan mennesker søker informasjon og
tar beslutninger. Designet og gjennomførte egne eksperimenter og undersøkelser,
og analyserte dem økonometrisk. Underviste og veiledet i statistikk og metode.
}

\cvitem{10/2014 -- 08/2016}{\textbf{Forskningsassistent} \newline
\textit{Istanbul Universitet -- Fakultet for økonomi}}
\cvitem{}{\vspace{-4pt}
Datainnsamling og kvantitativ analyse i prosjekter innen offentlig økonomi.
}

%----------------------------------------------------------------------------------------
%   UTDANNING
%----------------------------------------------------------------------------------------

\section{Utdanning}

\cvitem{2016 -- 2021}{\textbf{PhD i økonomi} \newline
\textit{Norges Handelshøyskole (NHH) -- FAIR, The Choice Lab} \newline
Avhandling: \textit{Essays on Information Preferences and Morality} \newline
Veiledere: Erik Ø. Sørensen (NHH) og Pedro Dal Bó (Brown University) \newline
Forskningsopphold ved Brown University (2019)}

\cvitem{2014 -- 2016}{\textbf{Mastergrad i økonomi} \newline
\textit{Galatasaray Universitet, Tyrkia}}

\cvitem{2009 -- 2014}{\textbf{Bachelor i økonomi og offentlig finans} \newline
\textit{Hacettepe Universitet, Tyrkia}}

%----------------------------------------------------------------------------------------
%   PUBLIKASJONER
%----------------------------------------------------------------------------------------

\section{Publikasjoner og analyser (utvalg)}

\cvitem{2026}{Hetland, A. G., Ay, F. C., Serdarevic, N., \& Lokken, N. C.
\textit{Noen Airbnb-utleiere rapporterer ikke leieinntektene sine, hvorfor?}
Skatteetatens analysenytt.
\href{https://www.skatteetaten.no/om-skatteetaten/analyse-og-rapporter/analysenytt/airbnb-utleiere-rapporter-ikke-leieinntektene-sine/}{\color{blue}\mbox{skatteetaten.no}}}

\cvitem{2025}{Ay, F. C., Freddi, E., Følstad, A., Hodnebrog, S., Kvale, K., Sell, O. A., \& Ulsaker, S.
\textit{Conversation logs as a source of insight: Predicting user satisfaction for customer service chatbots.}
Quality and User Experience, 10(1).}

\cvitem{2024}{Ay, F. C., Berge, J. W., \& Nødtvedt, K. B.
\textit{Strategic curiosity: An experimental study of curiosity and dishonesty.}
Journal of Economic Behavior \& Organization, 217, 287--297.}

\cvitem{2023}{Azevedo, F., Pavlovic, T., Rigo, G. G., Ay, F. C., \textit{et al.}
\textit{Social and moral psychology of COVID-19 across 69 countries.}
Nature Scientific Data, 10(1), 272.}

\cvitem{2022}{Van Bavel, J. J., Cichocka, A., Capraro, V., \textit{et al.} (inkl. Ay, F. C.)
\textit{National identity predicts public health support during a global pandemic.}
Nature Communications, 13(1), 517.}

\cvitem{}{Fullstendig publikasjonsliste: \href{https://cerenay.github.io}{\color{blue}cerenay.github.io}}

\end{document}
